\section{An explicit secant plane in every dimension}\label{sec:construction}

\Cref{q:higher} asks for four things: an abelian $n$-fold $X$, a plane
$P\subset H^{\mathrm{ev}}(X,\QQ)$ of Hodge classes meeting the spinor variety
in two points conjugate over $K$, a pair of secant sheaves realising the
plane, and a deformation argument. This section supplies the first three, for
every $n\ge2$ and every imaginary quadratic field $K$. The first two are
written in closed form (\Cref{thm:explicitweil,thm:plane}). The third is first
reduced to an explicit equation on a Chern character
(\Cref{thm:secantinvariants}) and then solved, in every dimension, by sheaves
of positive rank whose classes are combinations of the powers of a single line
bundle (\Cref{thm:secantexists}). The fourth ingredient is the deformation
argument, and \Cref{rem:deformation} says exactly what it requires.

The framework is Markman's \cite{Mar25a,Mar25b}: the Weil structure on
$X\times\widehat{X}$ obtained by embedding $K$ into
$\End^{0}(X\times\widehat{X})$, the notion of a $K$-secant sheaf, whose Chern
character lies in a plane meeting the spinor variety at two conjugate points,
and the deformation step that completes the argument at $n=3$ are all his.
What we add is the explicit form of the first three ingredients, uniformly in
$n$ and $K$. The closed form rests on two facts: the two conjugate isotropic
subspaces demanded by \Cref{cor:orbitcriterion} are the graphs of multiples of
a single element of the N\'eron--Severi group, and the pure spinor of the graph
of a two-form is the exponential of that form.

\subsection{The Weil structure}

\begin{setupenv}\label{setup:construction}
Let $X$ be a complex abelian variety of dimension $n$ with
$H_{1}(X,\QQ)=U$ of dimension $2n$ and complex structure $J$ on
$U\otimes\RR$. Let $d$ be a squarefree positive integer, $s=\sqrt{-d}$ and
$K=\QQ(s)$. Let $\beta\in\NS(X)\otimes\QQ$ be nondegenerate, regarded as an
alternating form on $U$, so that the Riemann relation $J^{\top}\beta J=\beta$
holds, and write $\phi_{\beta}\colon X\to\widehat{X}$ for the associated
isogeny. Put
\[
  A=X\times\widehat{X},\qquad \dim A=2n,\qquad
  H_{1}(A,\QQ)=U\oplus U^{\vee},
\]
and let
\begin{equation}\label{eq:Mdef}
  M=\begin{pmatrix}0 & -\beta^{-1}\\[2pt] d\,\beta & 0\end{pmatrix}
  \in\End^{0}(A),
\end{equation}
acting on $H_{1}(A,\QQ)=U\oplus U^{\vee}$ by
$(u,\phi)\mapsto(-\beta^{-1}\phi,\;d\,\beta u)$.
\end{setupenv}

\begin{theorem}[An explicit Weil structure]\label{thm:explicitweil}
In \Cref{setup:construction}:
\begin{enumerate}[label=\textup{(\roman*)}]
\item $M^{2}=-d$, so $\iota(s)=M$ defines an embedding
  $\iota\colon K\hookrightarrow\End^{0}(A)$;
\item $M$ commutes with the complex structure of $A$, and for an
  alternating form $\beta$ on $U$ this is equivalent to the Riemann relation,
  that is to $\beta\in\NS(X)\otimes\QQ$;
\item $M$ acts on $H^{1,0}(A)$ with eigenvalues $s$ and $-s$, each of
  multiplicity $n$, so $(A,\iota,\eta)$ is of $(K,-1,n)$-Weil type for a
  suitable polarisation $\eta$;
\item the eigenspaces $W_{\pm}\subset H_{1}(A,\CC)$ of $M$ for $\pm s$ are
  maximal isotropic for the canonical pairing on $U\oplus U^{\vee}$, are
  interchanged by complex conjugation, and are the graphs of the two-forms
  $\mp s\beta$.
\end{enumerate}
\end{theorem}

\begin{proof}
(i) Block multiplication gives
$M^{2}=\mathrm{diag}(-\beta^{-1}\,d\beta,\;-d\beta\,\beta^{-1})=-d\cdot\mathrm{id}$.

(ii) The complex structure of $A=X\times\widehat{X}$ is
$J_{A}=\mathrm{diag}(J,-J^{\top})$. The relation $MJ_{A}=J_{A}M$ amounts to
the two identities $d\beta J=-dJ^{\top}\beta$ and
$-\beta^{-1}(-J^{\top})=J(-\beta^{-1})$, and each is equivalent to
$\beta J=-J^{\top}\beta$. This is the Riemann relation: from
$J^{\top}\beta J=\beta$ and $J^{2}=-1$ one gets
$J^{\top}\beta=\beta J^{-1}=-\beta J$, as required. Conversely
$\beta J=-J^{\top}\beta$ gives $J^{\top}\beta J=-\beta J J=\beta$.

(iii) $M$ is block anti-diagonal and $H^{1,0}(A)=H^{1,0}(X)\oplus
H^{1,0}(\widehat{X})$ with both summands of dimension $n$, and $M$ exchanges
them. If $M(x,y)=\lambda(x,y)$ then $-\beta^{-1}y=\lambda x$ and
$d\beta x=\lambda y$, so $y=d\beta x/\lambda$ and, substituting,
$-dx/\lambda=\lambda x$, that is $\lambda^{2}=-d$. For each of
the two values of $\lambda$ the assignment $x\mapsto(x,d\beta x/\lambda)$ is an
isomorphism from $H^{1,0}(X)$ onto the $\lambda$-eigenspace, which is
therefore of dimension $n$. That is the balanced signature of
\Cref{def:weiltype}(ii). For (iii) of that definition, let $\eta_{0}$ be any
polarisation of $A$ and put
$\eta(x,y)=\eta_{0}(x,y)+d^{-1}\eta_{0}(Mx,My)$. A positive integer multiple
$mM$ of $M$ is an isogeny, so the second term is the positive multiple
$(dm^{2})^{-1}(mM)^{*}\eta_{0}$ of a polarisation, and $\eta$ is a
polarisation. Since $M^{2}=-d$,
$\eta(Mx,My)=\eta_{0}(Mx,My)+d\,\eta_{0}(x,y)=d\,\eta(x,y)$, and as
$M^{-1}=-M/d$ this gives $\eta(Mx,y)=-\eta(x,My)$: the Rosati involution of
$\eta$ sends $M$ to $-M$, that is $\iota(\tau)^{\dagger}=\iota(\bbar{\tau})$.
When $\beta$ is a polarisation an explicit choice of $\eta$ is recorded in
\Cref{prop:splitgeom}(i).

(iv) The same computation identifies $W_{\lambda}=\{(x,d\beta x/\lambda)\}$,
the graph of the form $d\beta/\lambda$, and
$d\beta/\lambda=d\beta/(\pm s)=\mp s\beta$ since $d/s=-s$. The canonical
pairing on $U\oplus U^{\vee}$ is
$\langle(x,\xi),(x',\xi')\rangle=\xi'(x)+\xi(x')$, and on the graph it gives
$d\beta(x',x)/\lambda+d\beta(x,x')/\lambda=0$ because $\beta$ is alternating; so
$W_{\lambda}$ is isotropic, and it has dimension $2n$, half of
$\dim(U\oplus U^{\vee})=4n$, hence is maximal. Conjugation sends $s$ to $-s$
and so interchanges the two.
\end{proof}

The placement of the factor $d$ in \eqref{eq:Mdef} is a convention, fixed once
and for all. The endomorphism $\begin{psmallmatrix}0&-d\beta^{-1}\\
\beta&0\end{psmallmatrix}$ also squares to $-d$, commutes with the complex
structure and defines a Weil structure on the same $A$, but it is a different
embedding of $K$: its eigenspaces are the graphs of $\mp s\beta/d$ rather
than of $\mp s\beta$, its polarisation is $\beta+d\widehat\beta$ rather than
$d\beta+\widehat\beta$, and the class in
$\bigwedge^{2}V_{+}\oplus\bigwedge^{2}V_{-}$ it produces is
$\beta-d\widehat\beta$ rather than $\gamma=d\beta-\widehat\beta$. The
convention of \eqref{eq:Mdef} is the one whose action on cohomology is
$x\mapsto-Bx$, $\xi\mapsto dB^{-1}\xi$ in the proof of
\Cref{thm:splitclosed}, and every formula below refers to it. For $d=1$ the
two conventions coincide.

\begin{remark}\label{rem:abundant}
\Cref{thm:explicitweil} says that abelian $2n$-folds of Weil type of the split
form $X\times\widehat{X}$ exist for every $n$, every imaginary quadratic $K$
and every polarised abelian $n$-fold $X$, and that the Weil structure is
written down by one element of $\NS(X)$. These are not general members of the
moduli of Weil type. The Weil classes on them are of interest because
Markman's deformation argument, where it can be carried out, propagates
algebraicity from them to the general member; see \Cref{rem:deformation}.
\end{remark}

\subsection{The secant plane}

Write $t=\beta\in\NS(X)\otimes\QQ$, so that by
\Cref{thm:explicitweil}(iv) the two conjugate maximal isotropic subspaces are
the graphs of $\mp s\,t$; we use a separate letter because everything below
holds for any nondegenerate $t\in\NS(X)\otimes\QQ$. The pure spinor
attached to the graph of a two-form $\omega$ is $e^{\omega}$, so the two
conjugate points of the even spinor variety determined by the Weil structure
are $[e^{-st}]$ and $[e^{st}]$.

\begin{definition}[The secant plane]\label{def:secantplane}
With $t$ and $d$ as above, put
\begin{equation}\label{eq:uv}
  u_{t}=\sum_{j\ge0}\frac{(-d)^{j}}{(2j)!}\,t^{2j},
  \qquad
  v_{t}=\sum_{j\ge0}\frac{(-d)^{j}}{(2j+1)!}\,t^{2j+1},
\end{equation}
both sums being finite because $t^{k}=0$ for $k>n$, and let
$P_{t}=\QQ u_{t}\oplus\QQ v_{t}\subset H^{\mathrm{ev}}(X,\QQ)$.
\end{definition}

\begin{theorem}[The plane is an explicit rational secant]\label{thm:plane}
With notation as above:
\begin{enumerate}[label=\textup{(\roman*)}]
\item $u_{t}$ and $v_{t}$ are rational Hodge classes, and
  $e^{\pm st}=u_{t}\pm s\,v_{t}$;
\item $P_{t}$ is a two-dimensional $\QQ$-subspace of $H^{\mathrm{ev}}(X,\QQ)$
  spanned by Hodge classes, and $\PP(P_{t})$ is a line defined over $\QQ$
  passing through the two conjugate points $[e^{\pm st}]$ of the even spinor
  variety;
\item that line is not contained in the spinor variety, so it meets it in
  exactly those two points.
\end{enumerate}
\end{theorem}

\begin{proof}
(i) Each $t^{k}$ is a rational Hodge class because $t\in\NS(X)\otimes\QQ$ is
of type $(1,1)$ and the algebraic classes form a subring. Separating
$e^{st}=\sum_{k}s^{k}t^{k}/k!$ into even and odd powers of $s$ and using
$s^{2j}=(-d)^{j}$, $s^{2j+1}=(-d)^{j}s$ gives $e^{st}=u_{t}+sv_{t}$, and
conjugating gives $e^{-st}=u_{t}-sv_{t}$.

(ii) The two elements are linearly independent over $\QQ$ because $u_{t}$ has
nonzero component in $H^{0}$ and $v_{t}$ does not, while $v_{t}\ne0$ since
$t\ne0$. From (i), both $e^{st}$ and $e^{-st}$ lie in $P_{t}\otimes K$, so
$\PP(P_{t})$ passes through the two points.

(iii) The even spinor variety is cut out by quadrics, by Cartan's description
of pure spinors, so a line either lies inside it or meets it in a subscheme of
length at most two. The point $[u_{t}]$ of $\PP(P_{t})$ is not a pure spinor:
a pure spinor with nonzero component in $H^{0}$ is a scalar multiple of
$e^{\omega}$ for a two-form $\omega$, and $u_{t}$ has zero component in
$H^{2}$ while its component in $H^{4}$ is $-\tfrac{d}{2}t^{2}\ne0$, whereas
$e^{\omega}$ with vanishing $H^{2}$ part is $1$. So the line is not contained
in the spinor variety. The two points of (ii) are distinct, since
$u_{t}+sv_{t}$ and $u_{t}-sv_{t}$ have the same component $1$ in $H^{0}$ and
$v_{t}\ne0$; hence they exhaust the intersection.
\end{proof}

\Cref{fig:spinor} draws the case $n=2$, in which the line $\PP(P_{t})$
meets the spinor quadric in the two conjugate points $[e^{\pm st}]$.

\begin{example}\label{ex:planes}
For $n=4$ and $d=1$,
$u_{t}=1-\tfrac{1}{2}t^{2}+\tfrac{1}{24}t^{4}$ and
$v_{t}=t-\tfrac{1}{6}t^{3}$. For $n=4$ and $d=3$,
$u_{t}=1-\tfrac{3}{2}t^{2}+\tfrac{3}{8}t^{4}$ and
$v_{t}=t-\tfrac{1}{2}t^{3}$. All four classes are rational and lie in the
subring $\QQ[t]$.
\end{example}

\subsection{The Chern character equation}

The next theorem computes the invariants that a sheaf with Chern character in
the secant plane must have. On the side of sheaves, the demands of
\Cref{q:higher} reduce to this one equation.

\begin{theorem}[The invariants of a secant sheaf]\label{thm:secantinvariants}
Let $X$, $t$, $d$ be as above and let $\cF$ be a coherent sheaf on $X$ with
$\ch(\cF)\in P_{t}$, say $\ch(\cF)=a\,u_{t}+b\,v_{t}$ with $a,b\in\QQ$. Then
\[
  \rk\cF=a,\qquad
  c_{1}(\cF)=b\,t,\qquad
  \ch_{2}(\cF)=-\tfrac{ad}{2}\,t^{2},\qquad
  c_{2}(\cF)=\tfrac{ad+b^{2}}{2}\,t^{2},
\]
and more generally $\ch_{k}(\cF)$ is the degree $k$ term of \eqref{eq:uv}
scaled by $a$ or $b$ according to the parity of $k$. The Bogomolov
discriminant is
\begin{equation}\label{eq:discnorm}
  \Delta(\cF)=2a\,c_{2}(\cF)-(a-1)c_{1}(\cF)^{2}
  =\Nm_{K/\QQ}\bigl(b+a s\bigr)\cdot t^{2}.
\end{equation}
In particular, if $t$ is ample and $\cF\ne0$, then
$\Delta(\cF)\cdot H^{n-2}>0$ for every ample $H$, so the Bogomolov inequality
places no obstruction on the existence of $\cF$.
\end{theorem}

\begin{proof}
The first four identities are read off from \eqref{eq:uv} degree by degree,
using $\ch_{2}=\tfrac12(c_{1}^{2}-2c_{2})$ to convert $\ch_{2}$ into $c_{2}$.
For \eqref{eq:discnorm},
\[
  2a\cdot\tfrac{ad+b^{2}}{2}t^{2}-(a-1)b^{2}t^{2}
  =\bigl(a^{2}d+ab^{2}-ab^{2}+b^{2}\bigr)t^{2}
  =\bigl(a^{2}d+b^{2}\bigr)t^{2},
\]
and $\Nm_{K/\QQ}(b+as)=(b+as)(b-as)=b^{2}+a^{2}d$. If $\cF\ne0$ then
$\ch(\cF)\ne0$, so $(a,b)\ne(0,0)$ and the norm is positive; and
$t^{2}H^{n-2}>0$ for $t$ and $H$ ample. This gives the last assertion.
\end{proof}

\begin{remark}[The norm form as the discriminant]\label{rem:normdisc}
Identity \eqref{eq:discnorm} says that the discriminant of the sheaf is the
norm form of $K$ evaluated at the coordinates $(a,b)$ of its Chern character
in the plane. The two conjugate points of the spinor variety are the two
isotropic directions of this norm form after extension of scalars to $K$, and
no nonzero rational point of the plane is isotropic. The Chern character of a
secant sheaf is a rational point of the line $\PP(P_{t})$, and so never an
isotropic one, while the Weil structure is carried by the two isotropic points
of that line, which are defined over $K$. In this sense the conjugate pair of
\Cref{cor:orbitcriterion} is built into the invariants of the sheaf.
\end{remark}

\begin{corollary}[The case of an abelian surface]\label{cor:surfacecase}
Let $n=2$, so $X$ is an abelian surface and $A=X\times\widehat{X}$ an abelian
fourfold of Weil type, and let $t$ be ample. A nonzero sheaf with
$\ch(\cF)\in P_{t}$ has Mukai vector
$v(\cF)=(a,\,bt,\,-\tfrac{ad}{2}t^{2})$ and Mukai self-pairing
\[
  \langle v(\cF),v(\cF)\rangle=\Nm_{K/\QQ}(b+as)\cdot t^{2}>0 .
\]
\end{corollary}

\begin{proof}
On an abelian surface the Mukai pairing of $v=(r,c,e)$ with itself is
$c^{2}-2re$, which here is $b^{2}t^{2}-2a(-\tfrac{ad}{2}t^{2})=(b^{2}+a^{2}d)t^{2}$.
It is positive because $(a,b)\ne(0,0)$ and $t^{2}>0$.
\end{proof}

\Cref{cor:surfacecase} is consistent with the case $n=2$ being settled: on an
abelian surface, positivity of the self-pairing of a primitive Mukai vector of
positive rank is sufficient for the moduli spaces of stable sheaves with that
vector to be non-empty.

\subsection{The genus three Jacobian}

At $n=3$ the equation of \Cref{thm:secantinvariants} determines a secant sheaf
of rank one almost completely, and it reproduces the shape of the object
Markman uses. The computation serves as a check on the framework of this
section.

\begin{proposition}[Rank one at $n=3$]\label{prop:genus3}
Let $X$ be the Jacobian of a smooth curve of genus three, let $\Theta$ be the
theta divisor and $C\subset X$ the Abel--Jacobi curve, so that
$\Theta^{2}/2=[C]$. Take $t=\Theta$ and $a=1$ in
\Cref{thm:secantinvariants}. Then a rank one torsion-free sheaf $\cF$ with
$\ch(\cF)\in P_{\Theta}$ and $c_{1}(\cF)=b\Theta$ is of the form
$\cF=\mathcal{I}_{Z}(b\Theta)$ with
\[
  [Z]=\frac{b^{2}+d}{2}\,\Theta^{2}=(b^{2}+d)\,[C].
\]
If $Z$ is a disjoint union of $m$ translates of $C$, then $b=1$ and
\[
  m=d+1 .
\]
\end{proposition}

\begin{proof}
A rank one torsion-free sheaf on a smooth projective variety is the tensor
product of its double dual, which is a line bundle, with the ideal sheaf of a
subscheme of codimension at least two. With $c_{1}=b\Theta$ this gives
$\cF=\mathcal{I}_{Z}(b\Theta)$ for such a subscheme $Z\subset X$, and
$\ch(\mathcal{I}_{Z}(b\Theta))=e^{b\Theta}\bigl(1-\ch(\cO_{Z})\bigr)$.
Comparing the components in $H^{4}$ with \Cref{thm:secantinvariants} at $a=1$,
which gives $\ch_{2}=-\tfrac{d}{2}\Theta^{2}$, yields
$\tfrac{b^{2}}{2}\Theta^{2}-[Z]=-\tfrac{d}{2}\Theta^{2}$, that is
$[Z]=\tfrac{b^{2}+d}{2}\Theta^{2}$. The Poincar\'e formula
$\Theta^{g-1}/(g-1)!=[C]$ at $g=3$ reads $\Theta^{2}/2=[C]$, whence
$[Z]=(b^{2}+d)[C]$.

Suppose now that $Z$ is a disjoint union of $m$ translates of $C$. Translates
of $C$ are homologous to $C$, so $m=b^{2}+d$. Each translate is a smooth curve
of genus three and $X$ has trivial tangent bundle, so Grothendieck--Riemann--Roch
gives $\ch(\cO_{Z})=m[C]-2m[\mathrm{pt}]$. Using $\Theta^{3}=6[\mathrm{pt}]$
and $\Theta\cdot[C]=3[\mathrm{pt}]$, the component of
$\ch(\mathcal{I}_{Z}(b\Theta))$ in $H^{6}$ is
$(b^{3}-3bm+2m)[\mathrm{pt}]$, and \Cref{thm:secantinvariants} requires it to
equal the degree three term $-bd\,\Theta^{3}/3!=-bd\,[\mathrm{pt}]$.
Substituting $m=b^{2}+d$ gives $-2(b-1)(b^{2}+d)=0$, so $b=1$ and $m=d+1$.
\end{proof}

\begin{table}[htbp]
\centering
\renewcommand{\arraystretch}{1.15}
\begin{tabular}{@{}rrrrrrr@{}}
\toprule
$d$ & $1$ & $2$ & $3$ & $5$ & $7$ & $11$ \\
\midrule
$m=d+1$ & $2$ & $3$ & $4$ & $6$ & $8$ & $12$ \\
\bottomrule
\end{tabular}
\caption{The number $m=d+1$ of translates of the Abel--Jacobi curve in
\Cref{prop:genus3}, as a function of the field $K=\QQ(\sqrt{-d})$.}
\label{tab:translates}
\end{table}

\begin{remark}\label{rem:markmancheck}
The sheaf used in \cite{Mar25b} on a genus three Jacobian is
$\mathcal{I}_{\cup C_{i}}(\Theta)$, the twist by the theta bundle of the ideal
sheaf of a union of translates of the Abel--Jacobi curve. This is the shape
that \Cref{prop:genus3} forces, including the twist by $\Theta$ itself, with
the number of translates determined by the field.
\end{remark}

\subsection{Existence of secant sheaves}\label{ssec:thirdexists}

\Cref{thm:secantinvariants} describes a secant sheaf but does not produce one.
We show that such sheaves exist on every polarised abelian variety, in
every dimension and for every imaginary quadratic field, and that their
classes can be assembled from the powers of a single line bundle.

Two observations do the work. The first is that the Chern character wanted is
a rational point of a plane, and a plane has no preferred scale: a positive
integer multiple $Mp$ of a rational point $p$ of $P_{t}$ is again a rational
point of $P_{t}$, and since the Hodge conjecture concerns classes with
rational coefficients, replacing $p$ by $Mp$ costs nothing. The second is
that, once denominators may be cleared, the Chern characters of the powers of
one line bundle span the subring $\QQ[t]$ in which the secant plane lies,
because the matrix expressing them in the basis $t^{k}/k!$ is a Vandermonde
matrix on distinct nodes. What remains is the passage from a class in $K_{0}$
to a sheaf, and for that the only hypothesis needed is that the rank be
positive.

\begin{lemma}[Positive rank is enough]\label{lem:posrank}
Let $Y$ be a smooth projective variety over a field and let $\alpha\in
K_{0}(Y)$ satisfy $\rk\alpha>0$. Then $\alpha=[\cF]$ for some coherent sheaf
$\cF$ on $Y$.
\end{lemma}

\begin{proof}
Write $m=\dim Y$ and $a=\rk\alpha\ge1$, and let $F^{c}K_{0}(Y)$ be the
topological filtration, the subgroup generated by the classes $[\cO_{Z}]$ of
closed integral subschemes $Z\subseteq Y$ of codimension at least $c$. Since
$Y$ is smooth, every coherent sheaf has a finite locally free resolution, so
$K_{0}(Y)=K_{0}'(Y)$ and the filtration is exhaustive with $F^{m+1}=0$.
The class of a coherent sheaf supported in codimension at least $c$
lies in $F^{c}$ and is congruent modulo $F^{c+1}$ to
$\sum_{i}\ell_{i}[\cO_{Z_{i}}]$, where $\sum_{i}\ell_{i}Z_{i}$ is its cycle in
codimension $c$; and $F^{1}$ is the kernel of the rank, while $F^{1}/F^{2}$ is
$\Pic(Y)$ through the first Chern class \cite[\S 15.1 and
Ex.~15.1.5]{Ful98}. We construct coherent
sheaves $\cF_{c}$ of rank $a$ with $\alpha-[\cF_{c}]\in F^{c}$, by induction on
$c$.

Start at $c=2$. Let $L$ be a line bundle with $c_{1}(L)=c_{1}(\alpha)$ and set
$\cF_{2}=L\oplus\cO_{Y}^{\oplus(a-1)}$. Then $\alpha-[\cF_{2}]$ has rank zero
and first Chern class zero, hence lies in $F^{2}$.

For the inductive step let $\delta=\alpha-[\cF_{c}]\in F^{c}$ with $c\ge2$.
Modulo $F^{c+1}$ the class $\delta$ is a finite sum $\sum_{j}m_{j}
[\cO_{Z_{j}}]$ with $m_{j}\in\ZZ$ and $\codim Z_{j}=c$. The terms with
$m_{j}>0$ are removed by replacing $\cF_{c}$ with $\cF_{c}\oplus\bigoplus_{
m_{j}>0}\cO_{Z_{j}}^{\oplus m_{j}}$, which changes neither the rank nor the
codimension of the remaining difference. For a term with $m_{j}<0$ we must
subtract, and here the rank is used. The restriction $\cF_{c}\otimes
\cO_{Z_{j}}$ is a quotient of $\cF_{c}$ and has rank at least $a\ge1$ at the
generic point of the integral scheme $Z_{j}$, by semicontinuity of the fibre
dimension. Dividing $\cF_{c}\otimes\cO_{Z_{j}}$ by a subsheaf of generic
corank one gives a surjection $\cF_{c}\twoheadrightarrow\cQ$ with $\cQ$ a
coherent $\cO_{Z_{j}}$-module of rank one. Then $[\cQ]\in F^{c}$ with
associated cycle $[Z_{j}]$, so $[\cQ]-[\cO_{Z_{j}}]\in F^{c+1}$. The kernel of
that surjection is a coherent sheaf, still of rank $a$ because $\cQ$ is
supported in positive codimension, and its class is $[\cF_{c}]-[\cQ]$.
Performing this $|m_{j}|$ times for each $j$ with $m_{j}<0$ produces
$\cF_{c+1}$ with $\alpha-[\cF_{c+1}]\in F^{c+1}$ and $\rk\cF_{c+1}=a$. At
$c=m+1$ the difference lies in $F^{m+1}=0$, so $\alpha=[\cF_{m+1}]$.
\end{proof}

\begin{theorem}[Secant sheaves exist in every dimension]\label{thm:secantexists}
Let $X$ be an abelian variety of dimension $n$, let $t\in\NS(X)\otimes\QQ$ be a
polarisation and let $d>0$. Let $p$ be a rational point of the secant plane
$P_{t}$ whose $u_{t}$ coordinate is positive. Then there are a positive
integer $M$, an integer $D>0$ with $Dt\in\NS(X)$, integers $n_{0},\dots,n_{n}$
and a coherent sheaf $\cF$ on $X$ such that, setting
\begin{equation}\label{eq:lbwitness}
  \alpha=\sum_{j=0}^{n}n_{j}\bigl[\cO_{X}(jDt)\bigr]\in K_{0}(X),
\end{equation}
one has $\alpha=[\cF]$ in $K_{0}(X)$, $\ch(\cF)=M\,p$ and
$\rk\cF=M\,a>0$, where $a$ is the $u_{t}$ coordinate of $p$. The integers $M$
and $n_{j}$ depend only on $n$, $d$, $D$ and the coordinates of $p$, and are
obtained by inverting one Vandermonde matrix. Moreover $Mp$ is again a
rational point of $P_{t}$; applying the statement to two linearly independent
rational points of $P_{t}$ produces a pair of coherent sheaves whose Chern
characters span $P_{t}$.
\end{theorem}

\begin{proof}
If the conclusion holds for $Np$, where $N$ is a positive integer, with
multiple $M'$, then it holds for $p$ with $M=NM'$, and the hypothesis on $p$
is unchanged by the scaling. Clearing the denominators of the coordinates of
$p$, we may therefore assume $p=au_{t}+bv_{t}$ with $a,b\in\ZZ$ and $a>0$.
Choose $D>0$ with $Dt\in\NS(X)$, so that $\cO_{X}(jDt)$ is a line bundle for
every $j\in\ZZ$.

Since $t$ is a polarisation, $t^{n}\ne0$ and $t^{n+1}=0$ in
$H^{\mathrm{ev}}(X,\QQ)$, so the $n+1$ classes $t^{k}/k!$ with $0\le k\le n$
are a basis of the subring $\QQ[t]$. In that basis
\[
  p=\sum_{k=0}^{n}c_{k}\,\frac{t^{k}}{k!},\qquad
  c_{k}=\begin{cases}a(-d)^{k/2}, & k\ \text{even},\\[2pt]
                     b(-d)^{(k-1)/2}, & k\ \text{odd},\end{cases}
\]
by \eqref{eq:uv}, and all the $c_{k}$ are integers. On the other
side,
\[
  \ch\bigl(\cO_{X}(jDt)\bigr)=e^{jDt}=\sum_{k=0}^{n}(jD)^{k}\,\frac{t^{k}}{k!},
\]
so in the same basis the $n+1$ vectors $\ch(\cO_{X}(jDt))$ for $j=0,\dots,n$
are the columns of the Vandermonde matrix $\bigl((jD)^{k}\bigr)_{0\le j,k\le n}$
on the $n+1$ distinct nodes $0,D,2D,\dots,nD$. Its determinant is
$\prod_{0\le i<j\le n}(jD-iD)=D^{n(n+1)/2}\prod_{0\le i<j\le n}(j-i)\ne0$, so
the system
\[
  \sum_{j=0}^{n}x_{j}\,(jD)^{k}=c_{k}\qquad(0\le k\le n)
\]
has a unique solution $x\in\QQ^{n+1}$. Let $M$ be the least common multiple of
the denominators of the $x_{j}$ and put $n_{j}=Mx_{j}\in\ZZ$. With $\alpha$ as
in \eqref{eq:lbwitness},
\[
  \ch(\alpha)=\sum_{j}n_{j}e^{jDt}
  =\sum_{k=0}^{n}\Bigl(\sum_{j}Mx_{j}(jD)^{k}\Bigr)\frac{t^{k}}{k!}
  =M\sum_{k=0}^{n}c_{k}\frac{t^{k}}{k!}=Mp .
\]
In particular $\rk\alpha=\ch_{0}(\alpha)=Mc_{0}=Ma>0$, so \Cref{lem:posrank}
gives a coherent sheaf $\cF$ with $[\cF]=\alpha$. The Chern character is a
function of the class in $K_{0}(X)$, whence $\ch(\cF)=Mp$ and $\rk\cF=Ma$.
Finally $Mp=(Ma)u_{t}+(Mb)v_{t}$ is a rational point of $P_{t}$, and two
linearly independent rational points stay linearly independent after scaling,
which gives the last assertion.
\end{proof}

\Cref{fig:moment} draws the theorem at $n=2$ and $d=2$, with the points of
the secant plane that integer combinations of line bundles reach.

\begin{remark}[The factorials and the multiple $M$]%
\label{rem:themultiple}
The integer $M$ is not an artefact of the proof. The lattice spanned by the
Chern characters of the powers of a single line bundle is spanned equally well
by the classes $(e^{t}-1)^{k}$ for $0\le k\le n$, because the $(n+1)$st finite
difference of a polynomial of degree $n$ vanishes; and $(e^{t}-1)^{k}=t^{k}+
\cdots$ has leading term $t^{k}$, not $t^{k}/k!$. The secant plane, by
contrast, is spanned by classes whose coordinates in the basis $t^{k}/k!$ are
integers. The index between the two lattices can be computed exactly. Write
$j^{k}=\sum_{i\le k}S(k,i)\,i!\binom{j}{i}$ with the Stirling numbers $S(k,i)$
of the second kind. As $j$ runs over $0,\dots,n$, the vectors
$\bigl(\binom{j}{i}\bigr)_{0\le i\le n}$ form a unitriangular basis of
$\ZZ^{n+1}$, so the vectors $(j^{k})_{0\le k\le n}$, which are the coordinates
of $\ch\cO_{X}(jt)$ in the basis $t^{k}/k!$, span the image of $\ZZ^{n+1}$
under the lower triangular matrix $T=\bigl(S(k,i)\,i!\bigr)_{k,i}$, whose
diagonal is $0!,1!,\dots,n!$. The index is therefore $\prod_{k\le n}k!$. The
inverse of $T$ is $\bigl(s(k,i)/k!\bigr)_{k,i}$, with the Stirling numbers
$s(k,i)$ of the first kind, because $\binom{j}{k}=\sum_{i}s(k,i)j^{i}/k!$;
so $n!\,T^{-1}$ is integral, and for every integral point $p$ of the secant
plane the least admissible $M$ divides $n!$. It is often larger than $1$. At $n=4$, $d=3$, $a=b=1$ the least multiple is $M=6$,
a divisor of $4!$, and the witness is
\[
  \alpha=-23\,[\cO_{X}]+66\,[\cO_{X}(t)]-54\,[\cO_{X}(2t)]
         +20\,[\cO_{X}(3t)]-3\,[\cO_{X}(4t)],
\]
of rank $-23+66-54+20-3=6=Ma$; in degree one its coefficient is
$66-108+60-12=6=Mb$, and in degrees two to four the sums
$\sum_{j}n_{j}j^{k}$ are $-18$, $-18$ and $54$, which are $6c_{2}$, $6c_{3}$ and
$6c_{4}$. The solution of the Vandermonde system is unique and the $n_{j}$
have no common factor, $23$ being prime to $6$, so $M=6$ is the least
multiple, and the freedom to rescale is needed. \Cref{fig:multiple} shows
the least multiple in further cases.

For a principal polarisation $\theta$ the classes $\theta^{k}/k!$ are
integral: in a symplectic basis $x_{1},\dots,x_{g},y_{1},\dots,y_{g}$ of
$H^{1}(X,\ZZ)$, $\theta=\sum_{i}x_{i}y_{i}$, the products $x_{i}y_{i}$
commute and square to zero, and so
$\theta^{k}/k!=\sum_{|I|=k}\prod_{i\in I}x_{i}y_{i}$. This is Poincar\'e's
formula in a form that does not need a curve, and it is why $t^{k}/k!$ is the
natural integral basis in which to take the coordinates $c_{k}$.
\end{remark}

\begin{figure}[htbp]
\centering
\includegraphics[width=0.85\textwidth]{fig_multiple}
\caption{\Cref{rem:themultiple}. The least multiple $M$ of
\Cref{thm:secantexists} for $1\le n\le10$, over the seventy-two cases
$d\in\{1,2,3,5,7,11\}$, $a\in\{1,2,3\}$, $b\in\{\pm1,\pm2\}$ for each
$n$, read off from the solution of the Vandermonde system, on a logarithmic
scale; the area of a disc counts the cases with that value. The solid line
joins the largest values and the dashed one is $n!$, which every $M$ divides
by the remark: the quotient of $\ZZ^{n+1}$ by the lattice of the
characters $(1,j,\dots,j^{n})$ of the powers of a line bundle is dual to the
integer-valued polynomials of degree at most $n$, spanned by the
$\binom{x}{k}$, whose coefficients have denominators dividing $n!$.}
\label{fig:multiple}
\end{figure}

\begin{figure}[htbp]
\centering
  \includegraphics[width=0.86\textwidth]{fig_moment}
\caption{\Cref{thm:secantexists} at $n=2$ and $d=2$. (a) The coordinates
$(\ch_{0},\ch_{1},\ch_{2})$ in the basis $1$, $t$, $t^{2}/2$. The Chern
characters of the powers of one line bundle lie on the moment curve
$j\mapsto(1,j,j^{2})$, and the Vandermonde determinant says that any three of
its points are affinely independent. The integral points of
$P_{t}=\QQ u_{t}\oplus\QQ v_{t}$ split by the criterion
$\ch_{2}\equiv\ch_{1}\bmod2$ into those reached by integer combinations of
line bundles, drawn filled, and those missed, drawn open. The point
$p=u_{t}+v_{t}$ is missed and $2p$ is reached, by the witness
$-3[\cO_{X}]+8[\cO_{X}(t)]-3[\cO_{X}(2t)]$, whose three terms are joined to
$2p$ by dashed segments. (b) The same points in the coordinates $(a,b)$ of
$a\,u_{t}+b\,v_{t}$. Here $\ch_{1}=b$ and $\ch_{2}=-2a$, so the criterion reads
$b\equiv0\bmod2$; the shaded half-plane $a>0$ is the range of
\Cref{thm:secantexists}.}
\label{fig:moment}
\end{figure}

\subsection{The three ingredients}

\begin{theorem}[Three of the four ingredients]\label{thm:twoofthefour}
For every $n\ge2$, every squarefree $d>0$ and every polarised abelian $n$-fold
$X$ carrying a nondegenerate $\beta\in\NS(X)\otimes\QQ$, the first three of the
four ingredients required by \Cref{q:higher} exist and are explicit:
\begin{enumerate}[label=\textup{(\alph*)}]
\item the abelian $2n$-fold of Weil type is $A=X\times\widehat{X}$ with the
  $K$-action \eqref{eq:Mdef}, by \Cref{thm:explicitweil};
\item the plane is $P_{t}$ of \Cref{def:secantplane} with $t=\beta$, it is
  spanned by Hodge classes and it meets the spinor variety in exactly the two
  conjugate points $[e^{\pm st}]$, by \Cref{thm:plane};
\item a pair of coherent sheaves $\cF_{1},\cF_{2}$ on $X$ whose Chern
  characters are linearly independent rational points of $P_{t}$ exists, by
  \Cref{thm:secantexists}, and each may be taken to have the class of an
  explicit integer combination of powers of one line bundle.
\end{enumerate}
The invariants of any such sheaf are forced by the explicit system of
\Cref{thm:secantinvariants}, which imposes $n-1$ linear conditions on a Chern
character already constrained to $\QQ[t]$, and, when $\beta$ is ample, are
unobstructed numerically by \eqref{eq:discnorm}.
\end{theorem}

\begin{proof}
Items (a), (b) and (c) are the cited theorems. For (c), the proof of
\Cref{thm:secantexists} uses of $t$ only that $t^{n}\ne0$ and that a positive
multiple of $t$ is integral; both hold for $t=\beta$, the first because
$\beta$ is nondegenerate. For the count in the last sentence, a Chern
character in the subring $\QQ[t]\subset H^{\mathrm{ev}}(X,\QQ)$ is a
polynomial of degree at most $n$ in $t$, so lies in a space of dimension
$n+1$, while $P_{t}$ has dimension two, and $(n+1)-2=n-1$.
\end{proof}

\begin{remark}[The deformation step]\label{rem:deformation}
The fourth ingredient of \Cref{q:higher} is untouched by this section, and it
is the hardest. Given the pair of secant sheaves $\cF_{1},\cF_{2}$ on $X$
supplied by \Cref{thm:secantexists}, one must still show that the
characteristic class of $\Phi(\cF_{1}^{\vee}\boxtimes\cF_{2})$ remains of
Hodge type under every deformation of $X\times\widehat{X}$ as a polarised
abelian $2n$-fold of Weil type, and that the object deforms with the variety,
as a twisted sheaf to first order. Markman carries out this step at $n=3$
\cite{Mar25b}, through a semiregularity analysis. Our contribution to it is
negative: \Cref{thm:factor} identifies the Hochschild classes that the
weakened form of that analysis must kill for a secant object, and
\Cref{thm:candidatefails} shows that the candidate Markman proposes for $n=4$
does not kill them. Without the step one obtains algebraicity only on the
split locus $X\times\widehat{X}$, where the Weil classes are of much less
interest, since the general abelian $2n$-fold of Weil type is simple and is
not of that form.

Accordingly, nothing in this section proves the Hodge conjecture for abelian
$2n$-folds of Weil type at any $n\ge4$. What is established is that three of
the four ingredients are written down for every $n$, the third with an
explicit witness in line bundles, and that the fourth is isolated.
\end{remark}

\begin{figure}[htbp]
\centering
  \includegraphics[width=0.88\textwidth]{fig_quadric}
\caption{The geometry of \Cref{thm:plane} in the first case, $n=2$. The even
spinor variety is then the smooth six-dimensional quadric of classes
isotropic for the Mukai pairing in $\PP(H^{\mathrm{ev}}(X,\CC))\cong\PP^{7}$,
and the figure shows a real two-dimensional section of it, a quadric surface
with its two rulings. The line $\PP(P_{t})$ of \Cref{def:secantplane} is
defined over $\QQ$ and meets the quadric only in the two points
$[e^{\pm st}]$, which are conjugate over $K$ and not individually rational;
they are drawn as real points. This is the configuration the Weil line
demands: a class defined over $\QQ$ cannot single out one of the two points,
so the construction must carry the pair, and the pair spans a plane rather
than a line. \Cref{prop:orbitplane} is the same statement on the
cohomological side, and \Cref{fig:escape} places it among the hypotheses a
construction can satisfy.}
\label{fig:spinor}
\end{figure}
